\begin{abstract}
El creciente volumen de datos geoespaciales generados por misiones satelitales y sistemas de exploración remota impone severas limitaciones en el ancho de banda disponible para transmisión y almacenamiento. Este artículo presenta un marco de compresión geoespacial basado en redes neuronales profundas que combina aprendizaje supervisado y self-supervised para lograr altas tasas de compresión manteniendo la utilidad analítica de los datos. La propuesta integra arquitecturas de compresión neuronal con módulos específicos para imágenes multiespectrales e hiperespectrales, evaluando tanto métricas rate-distortion tradicionales como el impacto en tareas downstream como clasificación de cobertura y detección de anomalías. Los experimentos demuestran reducciones significativas de bitrate (hasta 60-90% frente a VTM) con preservación de características críticas para exploración, superando enfoques convencionales y aprendidos genéricos. Se discuten implicaciones prácticas para misiones espaciales y limitaciones actuales.
\end{abstract}

\begin{IEEEkeywords}
Artificial Intelligence; Bandwidth Allocation; Deep Neural Networks (DNNs); Downstream Tasks; Earth Observation (EO); Hyperspectral Imaging; Image Compression; Multispectral Imaging; Neural Rate-Distortion; Remote Sensing; Self-Supervised Learning.
\end{IEEEkeywords}